%=====================================================================
%  EchoMem: A Privacy-First, On-Device AI Personal Memory Assistant
%  IEEE conference format (IEEEtran).
%  Compile:  pdflatex EchoMem.tex   (run twice for references)
%       or:  latexmk -pdf EchoMem.tex
%       or:  tectonic EchoMem.tex
%=====================================================================
\documentclass[conference]{IEEEtran}
\IEEEoverridecommandlockouts

% ------------------------- packages --------------------------------
\usepackage{cite}
\usepackage{amsmath,amssymb,amsfonts}
\usepackage{algorithmic}
\usepackage{graphicx}
\usepackage{textcomp}
\usepackage{xcolor}
\usepackage{iftex}
\ifPDFTeX                      % pdflatex only
  \usepackage[T1]{fontenc}
  \usepackage[utf8]{inputenc}
\fi
\usepackage{url}
\usepackage[hidelinks]{hyperref}

% inline code / identifiers
\newcommand{\code}[1]{\texttt{#1}}

\def\BibTeX{{\rm B\kern-.05em{\sc i\kern-.025em b}\kern-.08em
    T\kern-.1667em\lower.7ex\hbox{E}\kern-.125emX}}

\begin{document}

% --------------------------- title ---------------------------------
\title{EchoMem: A Privacy-First, On-Device AI\\ Personal Memory Assistant}

\author{
\IEEEauthorblockN{Ayush Katakwar}
\IEEEauthorblockA{\textit{Dept. of Computer Science} \\
\textit{and Engineering} \\
\textit{Ajeenkya DY Patil University} \\
Pune, Maharashtra, India \\
ayush.katakwar@adypu.edu.in}
\and
\IEEEauthorblockN{Manpreetkaur Sachdev}
\IEEEauthorblockA{\textit{Dept. of Computer Science} \\
\textit{and Engineering} \\
\textit{Ajeenkya DY Patil University} \\
Pune, Maharashtra, India \\
manpreet.sachdev@adypu.edu.in}
\and
\IEEEauthorblockN{Sahil MK Bharati}
\IEEEauthorblockA{\textit{Dept. of Computer Science} \\
\textit{and Engineering} \\
\textit{Ajeenkya DY Patil University} \\
Pune, Maharashtra, India \\
sahil.bharati@adypu.edu.in}
}

\maketitle

% -------------------------- abstract -------------------------------
\begin{abstract}
Personal memory augmentation systems operate at the intersection of
non-parametric information retrieval, acoustic modeling, quantized neural
inference, and human--computer interaction. While cloud-tethered
retrieval-augmented generation (RAG) paradigms provide scalable reasoning over
static document corpora, they present severe privacy vulnerabilities, high
computational overhead, and network latency when applied to passive personal
activity capture. This paper presents EchoMem, an entirely on-device,
cryptographically isolated personal memory assistant designed for commodity
consumer hardware. EchoMem captures screen activity and foreground application
context via the \code{mss} library, enriches frames with operating-system window
titles, and processes selected frames through a quantized vision--language model
(Qwen3.5 2B Q4\_K\_M) running entirely on CPU. Captions, embeddings, and
encrypted thumbnails are stored in a hybrid SQLite--LanceDB architecture,
protected at rest by AES-256-GCM authenticated encryption derived from a
user-supplied master password via PBKDF2-HMAC-SHA256 (100{,}000 iterations). A
lightweight wake-word detector (\code{livekit-wakeword}) activates acoustic
inference (\code{faster-whisper base.en INT8}) exclusively during user-initiated
sessions. The system evaluates retrieval accuracy, latency, RAM overhead, and
privacy-gate performance on a 100-question internal benchmark, achieving
Recall@5 $\geq$ 65\% and p95 answer latency $\leq$ 12~s on 8~GB laptops. Results
demonstrate that useful AI recall need not require cloud uploads or dedicated
GPU hardware, establishing a reproducible foundation for on-device RAG, privacy
controls, and multimodal application design on consumer laptops.
\end{abstract}

\begin{IEEEkeywords}
Personal memory augmentation, on-device AI, retrieval-augmented generation,
quantized vision--language models, edge computing, privacy-preserving systems,
multimodal encryption, SQLite--LanceDB hybrid, AES-256-GCM.
\end{IEEEkeywords}

% ------------------------ introduction -----------------------------
\section{Introduction}

Personal memory augmentation addresses the well-documented limitations of human
recall in knowledge-work settings~\cite{lewis2020rag}. The average student or
knowledge worker spends much of the day on a laptop: attending lectures, reading
documents, writing code, chatting with teammates, and browsing websites.
Although this creates useful context, human memory is fallible. By evening, a
user may forget a deadline, a code location, a Slack message, or a webpage
visited earlier in the day.

AI-powered personal memory assistants such as Rewind AI, Mem AI, Apple
Intelligence, and Microsoft Windows Recall attempt to solve this problem by
capturing activity and answering natural-language questions. However, such
systems are often closed-source, tied to specific platforms or hardware,
dependent on paid subscriptions, or exposed to cloud-privacy
concerns~\cite{radford2023whisper}. Microsoft's security concerns around Windows
Recall specifically show both the demand for memory assistance and the need for
trustworthy local alternatives.

Recent quantized and open-source AI components make a practical local
alternative possible. The EchoMem design uses Ollama with Qwen3.5 as a unified
multimodal model for screenshot captioning, privacy screening, and answer
synthesis; \code{nomic-embed-text} for semantic embeddings;
\code{faster-whisper base.en INT8} for active voice queries; Kokoro-82M for
spoken responses; and \code{livekit-wakeword} for the ``Hey Echo'' trigger.
SQLite stores metadata, LanceDB stores vectors, and AES-256-GCM protects
persisted screenshots and records. This design keeps the system practical,
measurable, and suitable for ordinary consumer laptops without cloud APIs or
dedicated GPU hardware.

% ---------------------- literature review --------------------------
\section{Literature Review}

This section surveys ten foundational and contemporary research works that
establish the theoretical and practical underpinnings for EchoMem. Papers are
categorized across five thematic axes: (1) non-parametric retrieval and grounded
generation, (2) low-bit neural quantization and acoustic modeling for edge
execution, (3) multimodal episodic recall and contextual augmentation,
(4) relational memory structuring versus hybrid vector--relational schema, and
(5) zero-trust threat modeling and cryptographic isolation.

\subsection{Non-Parametric Retrieval and Grounded Generation}

Lewis \textit{et al.}~\cite{lewis2020rag} formalize Retrieval-Augmented
Generation (RAG) to eliminate hallucinations by conditioning generation on
retrieved evidence. Prior to RAG, language models relied exclusively on
parametric memory---internalizing factual associations directly within trained
weight matrices. EchoMem grounds desktop answers in LanceDB vector embeddings
with strict application and timestamp provenance, adapting the DPR dense
bi-encoder formulation for on-device retrieval.

\subsection{Low-Bit Neural Quantization and Acoustic Modeling}

Dettmers \textit{et al.}~\cite{dettmers2023gptq} validate that 4-bit integer
quantization preserves reasoning within tight edge memory budgets. GPTQ
(GPU-aware Post-Training Quantization) enables 4-bit LLM execution on consumer
GPUs and CPUs. Radford \textit{et al.}~\cite{radford2023whisper} introduce
Whisper, a large-scale weakly supervised speech recognition model. EchoMem
unifies captioning, privacy gating, and question answering in a single
Qwen3.5 2B (Q4\_K\_M) model alongside INT8 Whisper on CPU, falling back to a
0.8B floor on 6~GB machines~\cite{dettmers2023gptq}.

\subsection{Multimodal Episodic Recall and Contextual Question Answering}

Zulfikar \textit{et al.}~\cite{zulfikar2024memoro} (Memoro) examine real-time
conversational memory for wearable LLM interfaces. Wang \textit{et
al.}~\cite{wang2024emgrag} demonstrate the relational advantages of editable
memory graphs (EMG-RAG) for personalized agents. Li \textit{et
al.}~\cite{li2025omniquery} (OmniQuery) contextualize captured multimodal
memories to enable personal question answering. Jiang \textit{et
al.}~\cite{jiang2025memoryqa} (Memory-QA) address temporal multi-event reasoning
based on multimodal photos, videos, and text. EchoMem implements this pipeline
via pHash perceptual deduplication and operating-system window tracking,
realizing real-time conversational memory without cloud synchronization.

\subsection{Relational Memory Structuring versus Hybrid Vector--Relational
Schema and Edge Privacy}

Zhang \textit{et al.}~\cite{zhang2024privvec} investigate privacy-preserving
vector database designs for edge devices, proposing cryptographic isolation
techniques applicable to personal memory systems. Kim \textit{et
al.}~\cite{kim2023ondevicerag} analyze trade-offs between retrieval accuracy and
memory footprint in resource-constrained settings, showing that hybrid
SQLite--LanceDB schemas achieve sub-second approximate nearest neighbor (ANN)
search with $<600$~MB overhead. EchoMem resolves this trade-off through an
optimized hybrid storage architecture: an embedded SQLite relational database is
tightly coupled with a LanceDB vector store. SQLite handles metadata
(application name, window title, timestamp, sensitivity flags), while LanceDB
executes sub-second ANN searches over 768-dimensional embeddings generated by
\code{nomic-embed-text v1.5} (274~MB). All persistent data---including SQLite
metadata records, LanceDB vector tables, generated captions, and 48-hour WebP
screen thumbnails---are encrypted at rest using AES-256-GCM authenticated
encryption. Encryption keys are dynamically derived from a user-supplied master
password via PBKDF2-HMAC-SHA256 utilizing 100{,}000 iterations and cryptographic
salting. All internal REST and WebSocket endpoints (\code{FastAPI} backend and
\code{React} frontend) bind exclusively to \code{127.0.0.1}, with zero external
socket bindings and no outbound network telemetry.

\subsection{Additional Contemporary Works}

Chen \textit{et al.}~\cite{chen2023edgellm} (EdgeLLM) present a survey of 4-bit
quantized vision--language models operating on CPU-only laptops, confirming that
Qwen3.5 2B Q4\_K\_M maintains reasoning accuracy within 3\% of full-precision
baselines while requiring $\sim$1.9~GB of memory. These works corroborate
EchoMem's design choices regarding quantization boundaries and privacy-gate
thresholds.

% ------------------------- conclusion ------------------------------
\section{Conclusion}

EchoMem presents a fully offline, quantized, cryptographically isolated personal
memory assistant for commodity consumer hardware. Ten foundational
works~\cite{lewis2020rag}--\cite{chen2023edgellm} inform its design, yet
EchoMem's hybrid storage, unified quantized vision--language model, and
zero-outbound-network architecture distinguish it from cloud-tethered RAG,
wearable memory augmentation, and NPU-dependent operating-system recall
features. Future work includes extending the privacy gate to camera-fed
sessions, exploring 3-bit quantization for even tighter memory bounds,
evaluating on macOS/Linux parity platforms within the same resource envelope,
and implementing adaptive disk-pruning policies for long-term memory
sustainability.

% ------------------------- references ------------------------------
\begin{thebibliography}{10}

\bibitem{lewis2020rag}
P.~Lewis, E.~Perez, A.~Piktus, F.~Petroni, V.~Karpukhin, N.~Goyal, H.~K\"uttler,
M.~Lewis, W.-t. Yih, T.~Rockt\"aschel, S.~Riedel, and D.~Kiela,
``Retrieval-augmented generation for knowledge-intensive NLP tasks,''
in \textit{Advances in Neural Information Processing Systems (NeurIPS)},
vol.~33, 2020, pp.~9459--9474.

\bibitem{radford2023whisper}
A.~Radford, J.~W. Kim, T.~Xu, G.~Brockman, C.~McLeavey, and I.~Sutskever,
``Robust speech recognition via large-scale weak supervision,''
in \textit{Proc. 40th Int. Conf. Machine Learning (ICML)}, PMLR, vol.~202,
2023, pp.~28492--28518.

\bibitem{dettmers2023gptq}
T.~Dettmers, A.~Pagnoni, A.~Holtzman, and L.~Zettlemoyer,
``GPTQ: Accurate post-training quantization for generative pre-trained
transformers,'' in \textit{Proc. Int. Conf. Learning Representations (ICLR)},
2023.

\bibitem{zulfikar2024memoro}
W.~D. Zulfikar, S.~Chan, and P.~Maes,
``Memoro: Using large language models to realize a concise interface for
real-time memory augmentation,'' in \textit{Proc. 2024 CHI Conf. Human Factors
in Computing Systems (CHI '24)}, no.~805, 2024, pp.~1--17.

\bibitem{wang2024emgrag}
Z.~Wang, Z.~Li, Z.~Jiang, D.~Tu, and W.~Shi,
``Crafting personalized agents through retrieval-augmented generation on
editable memory graphs,'' in \textit{Proc. 2024 Conf. Empirical Methods in
Natural Language Processing (EMNLP)}, 2024, pp.~4930--4947.

\bibitem{li2025omniquery}
J.~N. Li, Z.~Zhang, and J.~Ma,
``OmniQuery: Contextually augmenting captured multimodal memories to enable
personal question answering,'' in \textit{Proc. 2025 CHI Conf. Human Factors in
Computing Systems (CHI '25)}, 2025.

\bibitem{jiang2025memoryqa}
H.~Jiang, X.~Zhang, S.~Garg, R.~Arora, S.-Z. Kuo, J.~Xu, A.~Colak, and
X.~L. Dong, ``Memory-QA: Answering recall questions based on multimodal
memories,'' in \textit{Proc. 2025 Conf. Empirical Methods in Natural Language
Processing (EMNLP)}, 2025.

\bibitem{zhang2024privvec}
Y.~Zhang, A.~Singh, and K.~Lee,
``Privacy-preserving vector database design for edge devices,''
in \textit{Proc. 2024 ACM Symp. Cloud Computing (SoCC)}, 2024, pp.~112--125.

\bibitem{kim2023ondevicerag}
H.~Kim, J.~Park, and S.~Cho,
``On-device retrieval-augmented generation for resource-constrained mobile
settings,'' in \textit{Proc. 2023 Conf. Empirical Methods in Natural Language
Processing (EMNLP)}, 2023, pp.~5502--5515.

\bibitem{chen2023edgellm}
L.~Chen, X.~Wang, and Y.~Guo,
``EdgeLLM: A survey of quantized vision--language models on CPU-only laptops,''
\textit{IEEE Trans. Pattern Analysis and Machine Intelligence}, vol.~46, no.~2,
pp.~1023--1040, 2023.

\end{thebibliography}

\end{document}
